% Problem record op_6f1b7742387af9e7. This TeX file is derived from
% database/problems_json/op_6f1b7742387af9e7.json by scripts/sync-tex.mjs; edit the JSON record
% and rerun the script rather than editing this file.
\section{Distillation of every Wigner-negative depolarized Strange state}
\label{sec:op_6f1b7742387af9e7}

\paragraph{Problem.}
Can every Wigner-negative state in the depolarized qutrit Strange-state family be distilled arbitrarily accurately by catalyst-free stabilizer operations?

A stabilizer operation is constructed from stabilizer ancillas, Clifford unitaries, Pauli measurements, classical randomness and feedforward, and discarding systems. Heralding is permitted where stated. It does not mean an arbitrary stabilizer-preserving channel, and no magic catalyst is free. These distinctions affect conversion theorems \sourcecite{ref:6f1b-zurel26}{Zurel26}, \sourcecite{ref:6f1b-fang26}{Fang26}.

Define the qutrit Strange state and its depolarized family by

\begin{equation}
|S\rangle=\frac{|1\rangle-|2\rangle}{\sqrt2},\qquad
\rho_\epsilon=(1-\epsilon)|S\rangle\langle S|+\epsilon\frac{I_3}{3}.
\label{eq:6f1b-1}
\end{equation}

For every $0\leq\epsilon<3/4$, can copies of $\rho_\epsilon$ be distilled arbitrarily accurately to $|S\rangle$ using qutrit stabilizer operations?

More precisely, for every target error $\delta>0$, some finite number of copies must be converted with positive success probability to a state within trace distance $\delta$ of $|S\rangle\langle S|$, without a magic catalyst.

The endpoint $3/4$ is exact, not a guessed noise threshold. In the standard Gross discrete Wigner representation, with the phase-space origin chosen at the negative entry,

\begin{equation}
W_{\rho_\epsilon}(0)=\frac{4\epsilon-3}{9},\qquad
W_{\rho_\epsilon}(u)=\frac{3-\epsilon}{18}\quad(u\neq0).
\label{eq:6f1b-2}
\end{equation}

These expressions follow by mixing the Strange-state Wigner distribution with the uniform distribution. Hence the state is Wigner-negative exactly when $\epsilon<3/4$ \sourcecite{ref:6f1b-veitch12}{Veitch12}, \sourcecite{ref:6f1b-prakash20}{Prakash20}.

This one-parameter problem tests, but is not equivalent to the full resolution of, the broader conjecture that Wigner negativity suffices for odd-prime-dimensional magic-state universality \sourcecite{ref:6f1b-zurel26}{Zurel26}, \sourcecite{ref:6f1b-prakash26}{Prakash26}.

The displayed definitions, constraints, and target bounds are recorded in Eqs.~\eqref{eq:6f1b-1}, \eqref{eq:6f1b-2}.

\subsection*{Status}

Unsolved

\subsection*{Source}

This precise formulation is editor wording based on the unresolved direction and limitations documented in the cited primary literature \sourcecite{ref:6f1b-zurel26}{Zurel26}\sourcecite{ref:6f1b-veitch12}{Veitch12}; it is not presented as a verbatim conjecture of those authors.

\subsection*{Progress}

\begin{itemize}
  \item Necessity is known. Wigner-positive states remain Wigner-positive under stabilizer protocols, including conditioning on allowed measurement outcomes. Odd-dimensional nonstabilizer states with nonnegative Wigner functions therefore give genuine bound magic. Their existence does not answer the question about the Wigner-negative region \sourcecite{ref:6f1b-veitch12}{Veitch12}.

  \item 2020 construction. The 11-qutrit ternary Golay protocol distills the Strange state with cubic noise suppression. In the depolarizing convention used above, its threshold is approximately $\epsilon=0.387$, well below $0.75$ \sourcecite{ref:6f1b-prakash20}{Prakash20} (Sections 3 and 5).

  \item February 3, 2026 revision; 2026 journal publication. Prakash and Singhal searched a large class of qutrit stabilizer codes up to length 23. They found over 600 length-23 CSS examples with cubic suppression, but none surpassed the Golay threshold. This was a substantial search, not an exhaustive theorem over every possible protocol \sourcecite{ref:6f1b-prakash26}{Prakash26}.

  \item March 2026 update. The quadratic-residue-code analysis continues to identify the Golay code as the best known Strange-state threshold and leaves the Wigner-negativity conjecture open \sourcecite{ref:6f1b-zurel26}{Zurel26}.

  \item Status: open, including this simple radial test case. The presently unclosed interval should not be labeled a proven Wigner-negative bound-magic region.
\end{itemize}

\subsection*{References}

\begin{enumerate}[label={},leftmargin=4.8em,labelsep=0.5em]
  \footnotesize
  \item[\textup{[Zurel26]}]\label{ref:6f1b-zurel26}
  M. Zurel, S. Jana, and N. de Silva, \emph{High-threshold magic state distillation with quantum quadratic residue codes}, \href{https://arxiv.org/html/2603.18560v1}{arXiv:2603.18560v1}, March 19, 2026. Introduction explicitly states the open conjectures; Sections 4.2–4.3 give the new constructions.

  \item[\textup{[Fang26]}]\label{ref:6f1b-fang26}
  K. Fang and Z.-W. Liu, \emph{One-shot distillation with constant overhead using catalysts}, Nature Communications 17, 6010 (2026); \href{https://arxiv.org/html/2410.14547v3}{arXiv:2410.14547v3}, September 2, 2026. See Definitions 1–2 and Theorem 7.

  \item[\textup{[Veitch12]}]\label{ref:6f1b-veitch12}
  V. Veitch, C. Ferrie, D. Gross, and J. Emerson, \emph{Negative Quasi-Probability as a Resource for Quantum Computation}, New Journal of Physics 14, 113011 (2012); \href{https://arxiv.org/abs/1201.1256}{arXiv:1201.1256}. Establishes the Wigner-positive obstruction and bound magic.

  \item[\textup{[Prakash20]}]\label{ref:6f1b-prakash20}
  S. Prakash, \emph{Magic state distillation with the ternary Golay code}, \href{https://doi.org/10.1098/rspa.2020.0187}{Proceedings of the Royal Society A 476, 20200187 (2020)}; arXiv:2003.02717. See the depolarizing-noise convention and Strange-state threshold.

  \item[\textup{[Prakash26]}]\label{ref:6f1b-prakash26}
  S. Prakash and R. Singhal, \emph{A Search for High-Threshold Qutrit Magic State Distillation Routines}, \href{https://arxiv.org/html/2408.00436v2}{arXiv:2408.00436v2}, February 3, 2026; published as \emph{Search for high-threshold qutrit magic-state distillation routines}, Physical Review A 113, 042404 (2026). See the Introduction, code-search results, and conclusion.
\end{enumerate}

\subsection*{Comment}

Compared with the general mixed-state conjecture, this family offers a single real parameter, an exact geometric boundary, and a well-defined existing benchmark. It is a natural place to investigate whether coding constraints force a universal threshold gap or merely reflect limitations of existing constructions.

The appropriate negative result would rule out all adaptive stabilizer protocols, not just CSS projections or a bounded code-length search. A positive result along this ray would still leave other Wigner-negative qutrit directions to analyze.

\subsection*{Field}

Quantum Resource Theory

\subsection*{Topic}

Quantum magic; Resource conversion

\subsection*{ID}

\texttt{op\_6f1b7742387af9e7}
